\structuralelement{Список иллюстративного материала}

\begin{enumerate}
      \item Рисунок~\ref{fig:omni_scheme}~--- Кинематическая схема омни-платформы с колёсами Mecanum.
      \item Рисунок~\ref{fig:omni_principle}~--- Принцип действия колеса Mecanum и формирование результирующей реакции опоры.
      \item Рисунок~\ref{fig:house_scene_1}~--- Сцена симуляции квартиры (вид 1).
      \item Рисунок~\ref{fig:house_scene_2}~--- Сцена симуляции квартиры (вид 2).
      \item Рисунок~\ref{fig:rviz}~--- Рабочая сцена в RViz: карта помещения, costmap, лидар, путь и маркер цели.
      \item Рисунок~\ref{fig:exp_mode_run}~--- Временная диаграмма режимов сопровождения для отдельного запуска.
      \item Рисунок~\ref{fig:exp_vis_run}~--- Видимость цели $\nu(t)$ для отдельного запуска.
      \item Рисунок~\ref{fig:exp_safety_run}~--- Минимальное расстояние до препятствия по данным лидара с порогами $d_{\text{pot}}$ и $d_{\text{col}}$.
      \item Рисунок~\ref{fig:exp_smooth_run}~--- Норма приращения управления для отдельного запуска.
      \item Рисунок~\ref{fig:exp_all_modes}~--- Доли времени по режимам сопровождения по запускам.
      \item Рисунок~\ref{fig:exp_all_vis}~--- Доля времени видимости цели по запускам.
      \item Рисунок~\ref{fig:exp_all_safety}~--- Безопасность: число столкновений, заходов в зону риска и предотвращённых контактов по запускам.
      \item Рисунок~\ref{fig:exp_all_smooth}~--- Среднеквадратичная норма приращения управления по запускам.
      \item Таблица~\ref{tab:nodes}~--- Ключевые узлы программного комплекса.
\end{enumerate}
